% Part: computability
% Chapter: tm-computations
% Section: church-turing-thesis

\documentclass[../../../include/open-logic-section]{subfiles}

\begin{document}

\olfileid{tur}{mac}{ctt}
\olsection{د چرچ--ټيورينګ اصل}

ټيورينګ ماشينونه بايد د اغېزمنې کړنلارې د مفهوم کره بديل وي۔ ټيورينګ
فکر کاوه چې هر څوک چې هم د اغېزمنې کړنلارې او هم د ټيورينګ ماشين په
مفهوم پوه شي، دا شهود به ولري چې هر هغه کار چې د اغېزمنې کړنلارې په
وسيله کېدے شي، د ټيورينګ ماشين په وسيله هم کېدے شي۔ دا ادعا له دې
خبرې ملاتړ اخلي چې د اغېزمنې کړنلارې د مفهوم ټول نور وړانديز شوي کره
بديلونه په امتدادي ډول د ټيورينګ ماشين له مفهوم سره معادل راووځي---
يعنې د تابعو کټ مټ هماغه سټ محاسبه کولے شي۔ دې ادعا ته
\emph{د چرچ--ټيورينګ اصل} ويل کېږي۔

\begin{defn}[د چرچ--ټيورينګ اصل]
\emph{د چرچ--ټيورينګ اصل} وايي چې هر هغه څيز چې د اغېزمنې کړنلارې
په وسيله محاسبه کېدے شي، د ټيورينګ په معنا محاسبه کېدونکے دے۔
\end{defn}

د چرچ--ټيورينګ اصل په دوو لارو کارول کېږي۔ لومړے ډول کارول ئې د
سستۍ يوه پلمه ده۔ فرض کړئ د يوه څيز د محاسبې د اغېزمنې کړنلارې
توضيح لرو، مثلاً په ``کاذب کوډ'' کښې۔ بيا د چرچ--ټيورينګ اصل کارولے
شو، څو دا ادعا توجيه کړو چې هماغه تابع يو ټيورينګ ماشين محاسبه کوي،
ان که موږ هغه ماشين په عمل کښې نۀ وي جوړ کړے۔

د چرچ--ټيورينګ اصل بل استعمال له فلسفي پلوه لا په زړه پورې دے۔
ثابتېدے شي چې داسې تابعې شته چې ټيورينګ ماشينونه ئې نۀ شي محاسبه
کولے۔ له دې څخه د چرچ--ټيورينګ اصل په کارولو پايله اخيستلے شو چې
هغه تابع په هېڅ ډول کړنلارې سره په اغېزمنه توګه نۀ شي محاسبه کېدے۔
ځکه که داسې کړنلاره موجوده وے، د چرچ--ټيورينګ اصل له مخې به ورته يو
ټيورينګ ماشين هم موجود وے۔ نو که ثابت کړے شو چې هېڅ ټيورينګ ماشين
هغه نۀ شي محاسبه کولے، اغېزمنه کړنلاره هم ورته نشته۔ په ځانګړې توګه،
د چرچ--ټيورينګ اصل د دې ادعا دپاره کارول کېږي چې د درېدو په نوم
مسئله نۀ يوازې د ټيورينګ ماشينونو په وسيله نۀ شي حل کېدے، بلکې په
هېڅ ډول په اغېزمنه توګه نۀ شي حل کېدے۔

\end{document}
